\section{引言}

随着中国城市化进程的加快，大量农村人口迁移至城市，尤其是大中城市，导致城市人口迅速增长，进而带来了巨大的住房需求压力。
尽管近年来城镇新建住房比例显著提升，但历史上积累的自建房始终在存量住房中占较大比例，
造成存量的住房与流动的人口之间的矛盾 \cite{RN1}。与此同时，商品住房市场持续发展，但也伴随着高房价、高租金等问题，
这进一步加剧了中低收入群体，特别是新市民、青年人和基础服务人员在住房获取上的困难，
“青年群体住房难”现象已逐步转变成普遍性社会问题，严重影响我国民生建设的发展 \cite{RN2}。

在此背景下，保障性租赁住房作为国家住房保障体系的重要组成部分，
日益成为缓解大城市住房结构性矛盾、保障中低收入群体基本居住需求、推动人口有序流动与城市可持续发展的关键政策工具。
2021年7月，国务院办公厅印发《关于加快发展保障性租赁住房的意见》，明确提出将保障性租赁住房纳入住房体系顶层设计，
重点面向无房新市民、青年人群体供给小户型、低租金的住房产品。此后，各级政府陆续出台配套政策，
通过多种手段支持保障性租赁住房建设，包括土地供应、财政补贴、金融支持等。政策推动之下，
全国多个城市加快布局保障性租赁住房，初步构建了以政府为主导、多主体参与、市场化运作的建设模式。

然而，尽管取得了一定成效，我国保障性租赁住房的发展仍面临诸多现实困境。
一方面，地方政府在资金、土地等资源配置上的压力较大，政策实施存在较强的不平衡性和区域差异性；
另一方面，由于缺乏可持续的市场机制支持，企业参与热情不高，住房品质与管理服务难以保障 \cite{RN3} \cite{RN4}。
此外，政策协同机制不完善、租赁监管机制薄弱、租户权益保障缺失等问题，也在一定程度上制约了保障性租赁住房的健康发展。

因此，系统梳理我国保障性租赁住房的发展现状与现实困境，分析典型城市实践，提出切实可行的政策对策，
具有重要的理论意义与实践价值。不仅有助于深化对住房保障制度改革路径的认识，推动我国住房体系从“购房导向”向“租购并举”转型；
也能为地方政府和市场主体提供决策参考，提高资源配置效率与治理能力。


